\documentclass[12pt, a4paper]{article}
\usepackage[utf8]{inputenc}
\usepackage[T1]{fontenc}
\usepackage[english]{babel}
\usepackage{geometry}
\geometry{top=2.5cm, bottom=2.5cm, left=2cm, right=2cm}
\usepackage{parskip}
\usepackage{microtype}
\usepackage[default]{sourcesanspro}
\usepackage{xcolor}
\usepackage{fancyhdr}
\usepackage[most]{tcolorbox}
\usepackage[hidelinks]{hyperref}
\usepackage{amssymb, amsmath}

\newcommand{\Book}{Principles of Mathematical Analysis}
\newcommand{\Chapter}{The real and complex number systems}
\definecolor{StructureColor}{RGB}{40, 80, 120}

\pagestyle{fancy}
\fancyhf{}
\lhead{\color{StructureColor}\textbf{\Book}}
\rhead{\small\Chapter}
\cfoot{\thepage}
\renewcommand{\headrulewidth}{0.4pt}
\renewcommand{\headrule}{\hbox to\headwidth{\color{StructureColor}\leaders\hrule height \headrulewidth\hfill}}

\newtcolorbox{exercisebox}[1]{
    enhanced, colback=StructureColor!5!white, colframe=StructureColor,
    coltitle=white, fonttitle=\bfseries\sffamily,
    title={#1}, sharp corners, boxrule=0.5mm
}

\begin{document}
\thispagestyle{fancy}

\begin{exercisebox}{Exercise 3}
Prove the following statements:
\begin{itemize}
    \item[(a)] If $x \neq 0$ and $xy = xz$ then $y = z$.
    \item[(b)] If $x \neq 0$ and $xy = x$ then $y = 1$.
    \item[(c)] If $x \neq 0$ and $xy = 1$ then $y = 1/x$.
    \item[(d)] If $x \neq 0$ then $1/(1/x) = x$.
\end{itemize}
\end{exercisebox}

\textbf{Solution}

Let $x,y,z \in R$ with $x \neq 0$.

\textbf{Part (a)}

Suppose that $xy=xz$. Note that
\[
    y= y \cdot 1= y \cdot (\frac{1}{x}\cdot x)=\frac{1}{x}\cdot(y\cdot x) = \frac{1}{x}\cdot(z\cdot x) = z
\]

Therefore $y=z$. 

\textbf{Part (b)}

Suppose that $xy=x$. Note that $xy = x \cdot 1$, from Part (a) this implies that $y=1$. 

\textbf{Part (c)}

Suppose that $xy=1$. Note that $xy = x \cdot \dfrac{1}{x}$, from Part (a) this implies that $y=\dfrac{1}{x}$. 

\textbf{Part (d)}

Note that

\[
\dfrac{1}{\dfrac{1}{x}} = \dfrac{1}{\dfrac{1}{x}} \cdot 1 = \dfrac{1}{\dfrac{1}{x}} \cdot \dfrac{1}{x} \cdot x =x
\]

Therefore$1/(1/x) = x$. 

All statements were proved. $\blacksquare$
\end{document}